% ═══════════════════════════════════════════════════════════════════════
%  SITUATION D'APPRENTISSAGE — FICHE MANUSCRITE (copie élève)
%  3ème — Thème 2 : Configurations du plan — Leçon 2 : Propriété de Thalès
%  8 séances de 55 min — 8 Synthèses (S1 à S8) — Année scolaire 2026-2027
%
%  Style « cahier » : écriture manuscrite (police Patrick Hand, OFL),
%  encre bleue, titres des Synthèses au stylo rouge souligné, marge rouge,
%  lignes bleu pâle, filigrane du professeur.
%  NB : aucun mode mathématique — tout est composé dans la police
%  manuscrite (fractions en ligne, ex. AM/AB), comme sur une vraie copie.
% ═══════════════════════════════════════════════════════════════════════
\documentclass[12pt]{article}
\usepackage[a4paper,top=1.7cm,bottom=1.6cm,left=1.7cm,right=1.7cm]{geometry}
\usepackage{xcolor}
\usepackage{graphicx}
\usepackage{fontspec}
\setmainfont[Path=/fonts/]{PatrickHand-Regular.ttf}

% ── Couleurs « stylo » ─────────────────────────────────────────────────
\definecolor{encre}{RGB}{29,58,143}        % encre bleue
\definecolor{rougestylo}{RGB}{200,40,40}   % stylo rouge (titres)
\definecolor{ligne}{RGB}{168,190,214}      % lignes bleu pâle du cahier
\definecolor{grisfil}{RGB}{130,130,130}    % filigrane

\linespread{1.32}
\setlength{\parindent}{0pt}
\setlength{\parskip}{5pt}
\color{encre}

% ── Titre de Synthèse (stylo rouge, souligné) ──────────────────────────
\newcommand{\synth}[2]{\bigskip{\Large\color{rougestylo}\underline{\textbf{Synthèse #1}} \;\;{\normalsize (#2)}\par}\smallskip}
% ── Ligne bleu pâle du cahier ──────────────────────────────────────────
\newcommand{\lignecahier}{\medskip{\color{ligne}\rule{\linewidth}{0.4pt}}\medskip\par}
% ── Pointillés à compléter ─────────────────────────────────────────────
\newcommand{\pt}[1]{\textcolor{encre!55}{#1}}

% ── En-tête, marge rouge, filigrane, pied de page ──────────────────────
\newcommand{\filigrane}{%
  \rlap{\raisebox{11.5cm}[0pt][0pt]{\hbox to 0pt{\hskip 2.2cm
    \rotatebox{30}{\scalebox{3}{\textcolor{grisfil!45}{\bfseries KODJONE Kodjo}}}\hss}}}}
\newcommand{\margerouge}{%
  \rlap{\raisebox{0cm}[0pt][0pt]{\hbox to 0pt{\hskip -1.05cm
    \textcolor{rougestylo}{\vrule width 0.8pt height 26.6cm depth 0pt}}\hss}}}
\makeatletter
\def\ps@cahier{%
  \def\@oddhead{\filigrane\margerouge\hfill{\small SA — Propriété de Thalès}\hfill\mbox{}}%
  \def\@oddfoot{\hfill\textcolor{encre}{M. KODJONE Kodjo — Prof de Maths}\hfill
                \fbox{\small\arabic{page}}\hfill}%
}
\makeatother
\pagestyle{cahier}

\begin{document}
\large

% ═══════════════════════ EN-TÊTE DU CAHIER ═══════════════════════
\begin{center}
{\LARGE\color{rougestylo}\underline{\textbf{SITUATION D'APPRENTISSAGE}}}\\[6pt]
Mathématiques — Classe de \textbf{3ème} — Année scolaire \textbf{2026-2027}\\
Thème 2 : Configurations du plan — Leçon 2 : \textbf{Propriété de Thalès}\\
Durée : \textbf{8 séances de 55 min} — \textbf{8 Synthèses} (S1 à S8)\\[8pt]
\pt{Nom : .............................................. \; Prénoms : ..............................................}\\[2pt]
\pt{Classe : 3ème ........ \; Date : ................................}
\end{center}

\lignecahier

% ═══════════════════════ LA SITUATION ═══════════════════════
{\Large\color{rougestylo}\underline{\textbf{La situation}}}\par
\medskip
Pour aménager l'entrée du terrain de sport du \textbf{LYCÉE KPETA}, le maçon a
tracé deux allées rectilignes qui se coupent en A. Sur la première allée, il a
planté deux piquets B et M, avec AB = 9 m et AM = 6 m. Sur la deuxième allée,
il a planté deux piquets C et N, avec AC = 7,5 m et AN = 5 m. Il a ensuite
tendu une corde entre les piquets M et N.

\begin{center}
{\color{encre}%
\setlength{\unitlength}{1mm}%
\begin{picture}(46,36)
  \put(2,2){\rule{36mm}{0.5mm}}
  \put(20,28){\rotatebox[origin=l]{-55.30}{\rule{31.62mm}{0.5mm}}}
  \put(2,2){\rotatebox[origin=l]{55.30}{\rule{31.62mm}{0.5mm}}}
  \put(8,10.7){\rule{24mm}{0.5mm}}
  \put(19.4,27.4){\rule{1.3mm}{1.3mm}}
  \put(1.4,1.4){\rule{1.3mm}{1.3mm}}
  \put(37.4,1.4){\rule{1.3mm}{1.3mm}}
  \put(7.4,10.1){\rule{1.3mm}{1.3mm}}
  \put(31.4,10.1){\rule{1.3mm}{1.3mm}}
  \put(18.7,28.8){\large A}
  \put(0,-4){\large B}
  \put(38.2,-4){\large C}
  \put(3.2,11.4){\large M}
  \put(33.4,11.4){\large N}
  \put(16,-8.2){\large 9 m}
  \put(5.8,4.0){\large 6 m}
  \put(30.0,15.4){\large 5 m}
  \put(33.2,23.2){\large 7,5 m}
  \put(17.4,11.9){\large corde}
\end{picture}}
\end{center}

\textbf{Consigne générale :} chaque question de cette situation est appelée une
\textbf{Synthèse} (Synthèse 1, Synthèse 2, …, Synthèse 8). Chaque Synthèse sera
résolue à l'étape 4 de la séance : \textbf{d'abord individuellement (5 min),
puis en groupes (10 min)}. Rédige chaque solution dans ton cahier.

\lignecahier

% ═══════════════════════ SYNTHESE 1 ═══════════════════════
\synth{1}{Séance 1}
\textbf{1)} Citer les triplets de points alignés de la situation, puis repérer
les segments [MN] et [BC].\par
\textbf{2)} Quelle configuration géométrique reconnaît-on ? La nommer et la
décrire (où se trouvent les points M et N ?).

% ═══════════════════════ SYNTHESE 2 ═══════════════════════
\synth{2}{Séance 2}
\textbf{1)} Calculer le quotient AM/AB, puis le quotient AN/AC ; simplifier
chaque quotient.\par
\textbf{2)} Les points M et N sont-ils placés dans le même ordre que B et C à
partir de A ?\par
\textbf{3)} Que peut-on conclure pour les droites (MN) et (BC) ? Citer la
propriété utilisée.

% ═══════════════════════ SYNTHESE 3 ═══════════════════════
\synth{3}{Séance 3}
Le contrôleur mesure la corde : MN = 4 m.\par
Calculer la distance BC entre les piquets B et C. La rédaction doit comporter :
la propriété de Thalès écrite, le remplacement par les valeurs connues, puis le
produit en croix.

% ═══════════════════════ SYNTHESE 4 ═══════════════════════
\synth{4}{Séance 4}
\textbf{Deuxième partie de la situation :} pour l'autre accès du terrain, deux
chemins rectilignes se coupent en O. Des piquets A et B sont plantés sur le
premier chemin, avec OA = 4 m et OB = 10 m ; des piquets A' et B' sont plantés
sur le deuxième chemin, avec OA' = 6 m. On sait que (AA') // (BB').

\begin{center}
{\color{encre}%
\setlength{\unitlength}{1mm}%
\begin{picture}(52,28)
  \put(2,2){\rotatebox[origin=l]{33.69}{\rule{28.84mm}{0.5mm}}}
  \put(2,2){\rotatebox[origin=l]{19.98}{\rule{46.82mm}{0.5mm}}}
  \put(14,10){\rule{10mm}{0.5mm}}
  \put(26,18){\rule{20mm}{0.5mm}}
  \put(1.4,1.4){\rule{1.3mm}{1.3mm}}
  \put(13.4,9.4){\rule{1.3mm}{1.3mm}}
  \put(25.4,17.4){\rule{1.3mm}{1.3mm}}
  \put(23.4,9.4){\rule{1.3mm}{1.3mm}}
  \put(45.4,17.4){\rule{1.3mm}{1.3mm}}
  \put(0,-4){\large O}
  \put(13.0,11.6){\large A}
  \put(24.4,6.4){\large A'}
  \put(23.4,19.6){\large B}
  \put(45.6,19.6){\large B'}
\end{picture}}\par
\smallskip
{\color{encre}\large (OA = 4 m ; OB = 10 m ; OA' = 6 m)}
\end{center}

\textbf{1)} Quelle configuration reconnaît-on ? La nommer.\par
\textbf{2)} Calculer la distance OB'.

% ═══════════════════════ SYNTHESE 5 ═══════════════════════
\synth{5}{Séance 5}
Le maçon veut placer la porte P du terrain en partageant le côté [BC], qui
mesure BC = 6 m, dans la proportion 1 : 2.\par
\textbf{1)} Calculer les longueurs BP et PC.\par
\textbf{2)} Écrire le programme de construction qui permet de placer le point P
aux instruments (demi-droite, reports égaux, parallèle), puis faire la figure.

% ═══════════════════════ SYNTHESE 6 ═══════════════════════
\synth{6}{Séance 6}
Sur le second plan du terrain, on donne : OA = 4 m ; OB = 8 m ; OA' = 6 m ;
OB' = 12 m. Les points sont alignés dans le même ordre.\par
Le maçon affirme que les cheminements (AA') et (BB') sont parallèles. A-t-il
raison ? Justifier la réponse : ordre des points, calcul des deux quotients,
comparaison, conclusion en citant la propriété utilisée.

% ═══════════════════════ SYNTHESE 7 ═══════════════════════
\synth{7}{Séance 7}
Sur le plan de l'entrée, on veut une nouvelle configuration vérifiant
AM/AB = 1/3.\par
\textbf{1)} Écrire le programme de construction de cette configuration.\par
\textbf{2)} Réaliser la figure aux instruments.\par
\textbf{3)} Justifier que la configuration obtenue a bien la propriété
demandée.

% ═══════════════════════ SYNTHESE 8 ═══════════════════════
\synth{8}{Séance 8 — question de synthèse}
Le contrôleur mesure BC = 6 m et la corde MN = 4 m.\par
\textbf{1)} Vérifier par le calcul que MN/BC = AM/AB = AN/AC.\par
\textbf{2)} En déduire la position de la droite (MN) par rapport à la droite
(BC).\par
\textbf{3)} Que se passerait-il si la corde mesurait MN = 5 m ? Expliquer.

\lignecahier

\begin{center}
{\Large Bon courage ! Résous chaque Synthèse dans ton cahier.}
\end{center}

\end{document}
